The PHENIX detector system has finite energy and momentum resolution.
The limit in the resolution leads to the broadening of the signals that are usually described very well by the gaussian approximation.
This gaussian broadening can be estimated by performing MC of \Kstar with $\Gamma=0$ to prevent the width of a resonance interfering with the measurement.
For this analysis 10 million $K^{*0}(892) \rightarrow pi^\pm K^\mp$ events were generated for each dataset distributed uniformly within $0 < \varphi < 2\pi$, $\lvert \eta \rvert < 0.5$, and $\lvert z_{\text{VTX}} \rvert < 30$ from which 3, 3, 2, and 2 millions were generated uniformly within $0.4 < p_T < 0.9$, $0.9 < p_T < 2$, $2 < p_T < 5$, and $5 < p_T < 12$ GeV/$c$.
At first events were weighed to match $e^{-p_T}$ distribution, and after the estimation of invariant $p_T$ spectra events were weighted to match the obtained spectra.

Invariant mass distributions were obtained for DC-PC1 tracks for different $p_T$ from which the signals of \Kstar were approximated with 2 gausses - one for the signal of \Kstar and the other for background.
An example of such approximation is shown in Fig. \ref{fig:widthless_approximation_example}.
From these approximations widths of signals for different $p_T$ were obtained. 
These width vs $p_T$ distributions were approximated with 2nd order polynomial. Results are shown on Fig. \ref{fig:gaussian_broadening}. 
These approximatinos were used to estimate the gaussian broadening in real data for different $p_T$ bins.

\begin{figure}[h!]
	\centering
   \begin{subfigure}{0.32\textwidth}
      \includegraphics[width=\textwidth]{GaussianBroadening/Run14HeAu200/KStar892_1.0-1.2.pdf}
      \subcaption{}
   \end{subfigure}
   \hfill
   \begin{subfigure}{0.32\textwidth}
      \includegraphics[width=\textwidth]{GaussianBroadening/Run14HeAu200/KStar892_1.0-1.2.pdf}
      \subcaption{}
   \end{subfigure}
   \begin{subfigure}{0.32\textwidth}
      \includegraphics[width=\textwidth]{GaussianBroadening/Run14HeAu200/KStar892_1.0-1.2.pdf}
      \subcaption{}
   \end{subfigure}
   \caption[Examples of gaussian fit of $\Gamma = 0$ \Kstar in MC]{Examples of gaussian fit of $\Gamma = 0$ \Kstar in MC within $1 < p_T < 1.2$ GeV for Run14 \HeAu (a), Run15 \pAl (b), Run15 \pAu (c)}
	\label{fig:widthless_approximation_example}
\end{figure}

\begin{figure}[h!]
	\centering
   \begin{subfigure}{0.32\textwidth}
      \includegraphics[width=\textwidth]{GaussianBroadening/Run14HeAu200/KStar892_sigmas.pdf}
      \subcaption{}
   \end{subfigure}
   \hfill
   \begin{subfigure}{0.32\textwidth}
      \includegraphics[width=\textwidth]{GaussianBroadening/Run14HeAu200/KStar892_sigmas.pdf}
      \subcaption{}
   \end{subfigure}
   \begin{subfigure}{0.32\textwidth}
      \includegraphics[width=\textwidth]{GaussianBroadening/Run14HeAu200/KStar892_sigmas.pdf}
      \subcaption{}
   \end{subfigure}
   \caption[Gaussian broadening estimations of $\Gamma = 0$ \Kstar in MC]{Gaussian broadening estimations of $\Gamma = 0$ \Kstar in MC for different $p_T$ bins (red markers) and their approximations (dashed black line) for Run14 \HeAu (a), Run15 \pAl (b), Run15 \pAu (c)}
	\label{fig:widthless_approximation_example}
\end{figure}
